% \iffalse meta-comment
%%%%%%%%%%%%%%%%%%%%%%%%%%%%%%%%%%%%%%%%%%%%%%%%%
% lua-font-config.dtx
% Additions and changes Copyright (C) 2024-2026 Clea F. Rees.
% Code from skeleton.dtx Copyright (C) 2015-2024 Scott Pakin (see below).
%
% This work may be distributed and/or modified under the
% conditions of the LaTeX Project Public License, either version 1.3c
% of this license or (at your option) any later version.
% The latest version of this license is in
%   https://www.latex-project.org/lppl.txt
% and version 1.3c or later is part of all distributions of LaTeX
% version 2008-05-04 or later.
%
% This work has the LPPL maintenance status \muaintained'.
%
% The Current Maintainer of this work is Clea F. Rees.
%
% This work consists of all files listed in manifest.txt.
%
% The file lua-font-config.dtx is a derived work under the terms of the
% LPPL. It is based on version 2.4 of skeleton.dtx which is part of
% dtxtut by Scott Pakin. A copy of dtxtut, including the
% unmodified version of skeleton.dtx is available from
% https://www.ctan.org/pkg/dtxtut and released under the LPPL.
%%%%%%%%%%%%%%%%%%%%%%%%%%%%%%%%%%%%%%%%%%%%%%%%%
% \fi
%
% \iffalse
%<*driver>
\RequirePackage{svn-prov}
\ProvidesFileSVN{$Id: lua-font-config.dtx 12046 2026-09-18 16:15:13Z cfrees $}[v0.3 \revinfo][\filename: \filebase]
\DefineFileInfoSVN[lfc]
\GetFileInfoSVN*
% ref. ateb Max Chernoff: https://tex.stackexchange.com/a/723294/ 
\def\MakePrivateLetters{\makeatletter\ExplSyntaxOn\endlinechar13}
\documentclass[british,a4paper]{ltxdoc}
% ltnews42 / ateb David Carlisle: https://tex.stackexchange.com/a/754799/ (sylwad Ulrike Fischer rhywle ond alla i ddim canfod o hyd iddo fe)
\partokencontext=0
\usepackage{babel}
\usepackage{lua-font-config}
\usepackage{parskip}
\usepackage{verbatim}
\usepackage{graphicx}
\usepackage{xurl}
% sicrhau hyperindex=false: llwytho CYN bookmark
\usepackage{hypdoc}% ateb Ulrike Fischer: https://tex.stackexchange.com/a/695555/
\usepackage{bookmark}
\usepackage{fancyref}
\usepackage{csquotes}
\MakeAutoQuote{‘}{’}
\MakeAutoQuote*{“}{”}
\EnableCrossrefs
\CodelineIndex
\RecordChanges
%^^A \OnlyDescription
%^^A \ExplSyntaxOn
%^^A \makeatletter
\DoNotIndex{\verb,\ProvidesPackageSVN,\NeedsTeXFormat,\ProcessKeyOptions,\ExplSyntaxOn,\ExplSyntaxOff,\RequirePackage,\NeedsTeXFormat,\DefineFileInfoSVN,\GetInfoSVN,\@ifundefined,\@ifl@t@r\ExplLoaderFileDate,\endinput,\PackageError,\PackageWarning,\PackageInfo,\keys_define:nn,\keys_set:nn,\ProcessKeysOptionns,\author,\title,\maketitle,\date,\revinfo,\ProcessKeysOptions,\MessageBreak,\fmtversion,\documentclass,\ExplLoaderFileDate,\filebase,\fileversion,\fileinfo,\filebase:,\filedate,\documentclass,\@ifl@t@r,\begin,\end,\item,\IfFormatAtLeastTF}
%^^A \makeatother
%^^A \ExplSyntaxOff
\definecolor{strawberry}{rgb}{1.000,0.000,0.502}
\definecolor{blueberry}{rgb}{0.000,0.000,1.000}
\definecolor{moss}{rgb}{0.000,0.502,0.251}
\hypersetup{%
  colorlinks=true,
  citecolor={moss},
  extension=pdf,
  linkcolor={strawberry},
  linktocpage=true,
  pdfcreator={TeX},
  pdfproducer={pdfeTeX},
  urlcolor={blueberry}%
}
\NewDocElement[%
  idxtype=opt.,
  idxgroup=options,
  printtype=\textit{opt.},
]{Opt}{option}
\NewDocElement[%
  idxtype=fn.,
  idxgroup=functions,
  printtype=\textit{fn.},
  macrolike=true,
]{Fn}{function}
\NewDocElement[%
  idxtype=var.,
  idxgroup=variables,
  printtype=\textit{var.},
  macrolike=true,
]{Var}{variable}
\NewDocumentCommand \val { m }
{%
  {\ttfamily =\,\meta{#1}}%
}
\ExplSyntaxOn
\NewDocumentCommand \vals { m }
{
  {
    \ttfamily = \, 
    \clist_use:nn { #1 } { \textbar }
  }
} 
\ExplSyntaxOff
\setcounter{IndexColumns}{2}
\makeatletter
\def\@xobeysp{\leavevmode\penalty100\ }
\makeatother
\begin{document}
  \DocInput{\filename}
\end{document}
%</driver>
% \fi
% \title{\filebase}
% \author{Clea F. Rees}
% \date{\filedate}
% \maketitle
% \begin{abstract}
% An opentype font configuration package for Lua\LaTeXe{}.
% The aim is to provide an alternative to \pkg{fontspec} which is both faster
% and more powerful, by doing as much processing as possible in Lua.
% To avoid loading fonts unnecessarily, the data used is corrected during 
% compilation and the results cached for future use.
%
% \emph{This package is experimental.}
% \end{abstract}
% 
% \section{Requirements}\label{sec:require}
%
% The package requires Lua\LaTeXe.
% It will not work with any other engine or format.
%
% \section{Usage}\label{sec:pkg}
%
% \emph{The user interface is currently \textbf{extremely} basic.}
%
% All that is properly supported is the choice of font family for the standard
% families --- i.e.~roman/serif, sans and typewriter/mono --- and configuration
% of additional fonts, again by family.
% In addition to this, scaling and raw feature selection are supported.
% That is, most things are supported in the sense that they are possible, but
% the interface is rudimentary in the extreme.
%
% That said, the automated configuration is quite extensive.
% By default, the package tries to configure font families which are as rich
% as possible, insofar as the fonts support this. 
% This includes
% \begin{itemize}
%   \item multiple widths and shapes, insofar as \textsc{nfss} supports them;
%   \item italic, oblique, upright italic, reverse italic, reverse oblique;
%   \item small-caps, italic small-caps, oblique small-caps\footnote{%
%     This should work whether small-caps are in separate fonts (e.g.~Latin 
%     Modern) or implemented by the \texttt{smcp} feature (e.g.~\TeX{}-Gyre
%     Pagella).};
%   \item optical sizes.
% \end{itemize}
% Of course, this depends on the quality of the data the fonts provide.
% It is not always good or, even, vaguely accurate. 
% If things do not fit, an effort is made to accommodate them as distinct font
% families, so they are still available at document level.
%
% At present, this is a bit chaotic and a bit hard to use, as the 
% results are ‘hidden’ from you (in, of course, plain sight).
% Even so, you should typically need to do much less to access the fonts you
% want than you might otherwise expect. 
%
% The advantage of this is that if a font has several widths, you can easily
% adjust the one used for bold by simply redefining the relevant default.
% This is not very user friendly yet, but it is not terrible.
% 
%
% Simple usage looks like this:
% \iffalse
%<*verb>
% \fi
\begin{verbatim}
  \usepackage{lua-font-config}
  \fontconfig{%
    rm = texgyrepagella,
    sf = texgyreheros,
    tt = texgyrecursor,
  }
\end{verbatim}
% \iffalse
%</verb>
% \fi
% 
% If you want to configure font features, scale a family or use a non-default
% script etc., you can use:
% \iffalse
%<*verb>
% \fi
\begin{verbatim}
  \usepackage{lua-font-config}
  \fontconfig{%
    rm = {%
      name = texgyrepagella, 
      features = {script=ltn;+tlig;+pnum},
    },
    sf = {%
      name = texgyreheros,
      scale = 1.2,
      features = {script=cyrl;},
    },
    tt = {%
      name = texgyrecursor,
      features = {+zero;},
    },
  }
\end{verbatim}
% \iffalse
%</verb>
% \fi
% Do \textbf{not} specify small-capitals (\texttt{+smcp}) unless you want 
% \emph{all} fonts in the family to use that feature.
% Small-caps and italic/oblique small-caps will be configured 
% automatically if available.
%
% \section{Advanced Configuration}\label{sec:adv}
%
% Advanced configuration is possible, but the interface is currently non-
% existent and, despite not existing, will certainly change. 
%
% \iffalse
%<*verb>
% \fi
\begin{verbatim}
\directlua{
  lfc = require("lfc")
  local font_config = lfc.font_config
  font_config("venturis", {fea = "mode=node;script=dflt;+tlig;+pnum;+onum"})
}
\end{verbatim}
% \iffalse
%</verb>
% \fi
%
% You do not actually gain much by this, although it is a bit faster, since 
% key-value processing in \LaTeX{} is fairly slow. 
% However, you cannot actually do anything you cannot do otherwise.
% The only two configuration options are \texttt{fea} and \texttt{scale},
% although you could theoretically modify them for different families, if you
% lookup string will return several.
%
% What you gain is primarily the absence of interference on the \LaTeX{} side.
% Families will be configured for you, but you may then choose how to use 
% them; nothing will change by default.
% Not even macros will be defined to access them; if you want that, you must
% do it yourself.
% 
% So after the Lua above, you might do
% \iffalse
%<*verb>
% \fi
\begin{verbatim}
\renewcommand\rmdefault{venturisadf}
\renewcommand\sfdefault{venturissansadfcd}
\makeatletter
% TitlingNo3? (digits get eaten during lookup)
\DeclareRobustCommand\venturistitling{%
  \not@math@alphabet\venturistitling\relax
  \fontfamily{venturisadftitlingno}\selectfont
}
\DeclareTextFontCommand\textvtitle{\venturistitling}
\makeatother
\end{verbatim}
% \iffalse
%</verb>
% \fi
% Right now, the only way to know what is available is to examine the fonts
% or, better, the cached data (after the first run).
% You can also put \texttt{lfc_debug = true} before loading the module to 
% get (lots) of console output, including the database contents.
%
% That is, you can use the module directly, but that is not to say that the
% module is actually usable at present.
%
% Obviously, such shenanigans are not for the faint of heart or, indeed, for
% the stout-hearted-but-moderately-sane.
%
% 
% \changes{v0.0}{0000-00-00}{First public release.}
%  
%
% \MaybeStop{%
%   \PrintChanges
%   \PrintIndex
% }
% 
% \section{Implementation}
%
% You do not need to read the remainder of this document in order to install or use the package.
%
% Note that almost everything is done in Lua, so there is really very little to see here.
%
%<@@=lfc>
%<*sty>
%    \begin{macrocode}
\NeedsTeXFormat{LaTeX2e}[2021-11-15]
\RequirePackage{svn-prov}
\ProvidesPackageSVN[\filebase.sty]{$Id: lua-font-config.dtx 12046 2026-09-18 16:15:13Z cfrees $}[v0.0 \revinfo][\filebase: Font configuration for LuaLaTeX]
% this is needed only because l3doc overwrites \fileversion and \filedate
% neither ltxdoc nor doc does this
\DefineFileInfoSVN[lfc]
%    \end{macrocode}
% \iffalse
% ^^A Paid â defnyddio \GetFileInfoSVN*/\GetFileInfoSVN{} yn y fan hon!!
% \fi
%    \begin{macrocode}
%%%%%%%%%%%%%%%%%%%%%%%%%%%%%%%%%%%%%%%%%%%%%%%%%
%% copied verbatim, excepting format from Joseph Wright's' siunitx.sty under LPPL
\@ifundefined{directlua}{%
  \PackageError{lua-font-config}{Package requires LuaLaTeX}
  {To use lua-font-config, you must compile your document with LuaLaTeX.}%
}{}
\@ifundefined{ExplLoaderFileDate}{%
  \RequirePackage{expl3}%
}{}
%% almost verbatim from siunitx.sty
\@ifl@t@r\ExplLoaderFileDate{2022-02-24}{%
}{%
  \PackageError{lua-font-config}{Support package expl3 too old}
  {%
    You need to update your installation of the bundles 'l3kernel' and
    'l3packages'.\MessageBreak
    Loading~lua-font-config~will~abort!%
  }%
  \endinput
}
%%%%%%%%%%%%%%%%%%%%%%%%%%%%%%%%%%%%%%%%%%%%%%%%%
%%%%%%%%%%%%%%%%%%%%%%%%%%%%%%%%%%%%%%%%%%%%%%%%%
\ExplSyntaxOn
\prop_gput:NnV \g_msg_module_name_prop {@@} \lfc@base
\msg_new:nnnn {@@} {needs-luatex} {
  \msg_critical_text:n {@@} :: ~requires~LuaLaTeX~\msg_line_context:
} { \msg_see_documentation_text:n {@@} }
\sys_if_engine_luatex:F { \msg_critical:nn {@@} {needs-luatex} }
%    \end{macrocode}
% \begin{variable}{\l__@@_tmpa_tl}
%   Custom ephemeral token list.
%    \begin{macrocode}
\tl_new:N \l__@@_tmpa_tl
%    \end{macrocode}
% \end{variable}
%    \begin{macrocode}
%    \end{macrocode}
% Configure keys for default families.
%    \begin{macrocode}
\clist_map_inline:nn {rm, sf, tt} {
  \keys_define:nn {lfc} {
    #1 .code:n            = { 
      \keys_set_exclude_groups:nnnN {lfc/#1} {config} {##1} \l__@@_tmpa_tl
      \keys_set:nV {lfc/#1} \l__@@_tmpa_tl
    },
    #1 .value_required:n  = true,
    #1 .usage:n           = preamble,
  }
  \keys_define:nn {lfc/#1} {
    name      .code:n = { \use:c {__@@_set_#1:n} {##1} },
    features  .code:n = { \__@@_set_fam_fea:n {##1} },
    unknown   .code:n = {
      \exp_args:NnV \use:c {__@@_set_#1:n} \l_keys_key_str
    },
    name      .value_required:n   = true,
    features  .value_required:n   = true,
    unknown   .value_forbidden:n  = true,
    name      .usage:n            = preamble,
    features  .usage:n            = preamble,
    unknown   .usage:n            = preamble,
  }
}
%    \end{macrocode}
% Configure keys for additional families.
%    \begin{macrocode}
\keys_define:nn {lfc} {
  family  .code:n           = { 
    \keys_set_exclude_groups:nnnN {lfc/fam} {config} {#1} \l__@@_tmpa_tl
    \keys_set:nV {lfc/fam} \l__@@_tmpa_tl
  },
  family  .value_required:n = true,
  family  .usage:n          = general,
}
\keys_define:nn {lfc/fam} {
  name      .code:n = { \__@@_set_fam_name:n {#1} },
  features  .code:n = { \__@@_set_fam_fea:n {#1} },
  unknown   .code:n = {
    \exp_args:NV \__@@_set_fam_name:n \l_keys_key_str 
  },
  name      .value_required:n   = true,
  features  .value_required:n   = true,
  unknown   .value_forbidden:n  = true,
  name      .usage:n            = general,
  features  .usage:n            = general,
  unknown   .usage:n            = general,
  features  .groups:n           = config,
}
%    \end{macrocode}
% \begin{macro}{\fontconfig}
%   The main (and currently(?) only) macro.
%    \begin{macrocode}
\NewDocumentCommand\fontconfig { +m }{
  \keys_set:nn {lfc} {#1}
  \__@@_configure_doc_families:
}
%    \end{macrocode}
% \end{macro}
%    \begin{macrocode}
\directlua{
  require("lfc")
}
\ExplSyntaxOff
%    \end{macrocode}
%</sty>
% 
%\Finale
